\section{Verification methodology}
\label{sec:proof-design}

The proof was organized around the data actually received and produced by the
verifier.  This choice shaped both the Lean model and the translation work.

\subsection{Prove local refinements, then one global execution}

For each verifier stage we use three objects: the Rust computation, its
functional translation, and a mathematical specification.  Local refinement
theorems connect the translated function to its specification.  A final
composition theorem starts from an arbitrary successful top-level call and threads the
same execution record through every local theorem.

This last step is essential.  Existential statements of the form ``some
accepted transcript has these challenges'' and ``some accepted opening has
these values'' can both be true while referring to different runs.  The
same-execution theorem instead proves
\[
  \rho\text{ succeeds}
  \Longrightarrow
  \bigwedge_k \mathcal P_k(\rho),
\]
where each property \(\mathcal P_k\) is projected from the same \(\rho\).
Intermediate equalities are obtained by reducing the translated code, not by
adding them to the theorem's hypotheses.

\subsection{Begin with proof bytes}

The chain verifier receives roots, selected values, salts, and compressed
authentication paths.  It never receives complete committed tables.  The
formal opening model therefore begins with the literal byte stream.  A
successful section proof must account for every value, salt, path slot, and
frontier node, and must return its exact unused suffix.  The five-section
routine passes those suffixes in order and accepts only an empty final suffix.

This makes omissions, duplicate records, wrong tree order, and trailing bytes
ordinary failed proof obligations.  Starting from a hypothetical full table
would hide all four cases behind an unproved parser assumption.

\subsection{Preserve causality}

Every prover message is indexed by the transcript prefix available when the
message is chosen.  Each bad-challenge set is likewise fixed before its
challenge is sampled.  This discipline is used in the four-fold extraction,
the relation repair bound, and the multi-attempt oracle argument.

The same rule applies at byte level.  The transcript model records absorbs,
squeezes, and work checks as distinct operations.  The nonce is checked
against the pre-absorb state, then absorbed, and only then may the next
challenge be produced.  Exact order is part of the theorem rather than a
comment on the implementation.

\subsection{Join authentication and algebra explicitly}

Merkle authentication proves that a value belongs to a commitment.
Low-degree testing proves an algebraic property of queried values.  Neither
fact implies the other.  A consumed-value map connects them by proving that
each value read by the fold and relation code is the value returned by the
corresponding authenticated opening in the same run.

Public-statement binding receives the same treatment.  The relation theorem
uses nineteen columns in a fixed order, and the public-field rows identify six
of those values with the spend statement.  Lane batching, relation folding,
and public binding are separate lemmas joined only at the final theorem.

\subsection{Extract collision witnesses}

SHA-256 and Poseidon2 compress their inputs, so their formal interfaces do not
assume injectivity.  When two reached inputs have the same output, the proof
either establishes equality of the inputs or returns the unequal pair as a
collision witness.  The security experiment can then charge the collision
event under a concrete primitive assumption.

Merkle paths add a positional condition.  Comparing two different leaves at
the same position identifies the first unequal ordered node inputs with one
equal parent.  Paths at different positions remain separate cases in the
theft game.

\subsection{Keep probability units visible}

Every false-acceptance event lives in one probability space, and the detailed
ledger has a proved union equality with its grouped form.  This prevents a
conditioning change from being hidden by prose.  It also forces the analysis
to distinguish work-normalized probability from the probability of failure
after the work search has finished.  The factor \(2^{37}\) belongs to one of
those experiments, not both.

\subsection{Audit the proof as an artifact}

The replay procedure resolves the dependency closure of the final theorem
from a clean checkout.  It rejects admitted proofs, unchecked native
evaluation, and generated files absent from the manifest.  It then compiles
the closure and prints the final axiom report.  Translation inputs, generated
code, build tools, deployed bytes, proof bytes, and transaction records are
checked by separate manifests.

This produces four kinds of evidence with different meanings:

\begin{description}[leftmargin=0.23\textwidth,style=nextline]
  \item[Machine-checked theorem] Lean accepts a statement under its displayed
        hypotheses.
  \item[Conditional reduction] The conclusion follows when a published
        mathematical or cryptographic assumption is supplied.
  \item[Deterministic identity] Recorded source, generated code, or program
        bytes match by replay and hash.
  \item[Execution record] One archived transaction exhibits the stated input,
        cost, and state change.
\end{description}

The paper combines these forms of evidence without treating one as a
substitute for another.
